\section{Adaptive Mesh Refinement (AMR)}

Loci provides Adaptive Mesh Refinement (AMR) features that can be
exploited by CHEM.  AMR features can be executed in either an offline
or an online mode.  The offline AMR features are accomplished by
executing stand-alone commands that can edit the refinement plans for
a given mesh specified as a VOlume Grid (vog) file.  The AMR technique
is designed to be able to refine any polyhedral mesh and refines a
cell through a technique known as midpoint insertion.  A cell is
refined by first splitting each edge of the cell at the edge midpoint.
Then the face centroid is inserted and a face is formed for each
vertex of the original face by adding edges between the edge midpoint
and the face centroid.  This process proceeds recursively to refine
the cell through inserting the cell centroid where every vertex of the
cell is formed into a daughter cell.  This refinement approach has the
advantage that it works on general meshes composed of arbitrary
polyhedral cells.


\subsection{{\tt adpt}: Generating the Tag File}
The \verb!adpt! code reads the Chem grid and solution files and tags
the nodes for refinement based on the error indicators given.  The
result of running \verb!adpt! is the tag file (named
\verb!refine.dat!) which consists of {\tt 1}'s and {\tt 0}'s to
indicate whether or not a node is attached to a cell that needs to be
refined.

The format for the \verb!adpt! command is:

\begin{verbatim}
adpt {opt args} <iter> <sensitivity> case <errorvars>
\end{verbatim}

\noindent
where

\begin{description}
\item {\tt \verb!{!opt args\verb!}!} :  optional arguments to reduce the extent of the refinement region (see below)
\item {\tt <iter>} :  iteration number to use for the tag file
\item {\tt <sensitivity>} :  error sensitivity (use a small number to tag more nodes for refinement)
\item {\tt case} :  case name
\item {\tt <errorvars>} :  list of error indicators to use
\end{description}
\noindent
Note that when using the script, the resulting \verb!refine.dat! file
is located in the level directory associated with the current
refinement level ({\it{e.g.}}, \verb!level.2/!).

The optional arguments allow the definition of a smaller region for
refinement through the specification of min and max coordinate values
and/or min and max distances to viscous walls.  If the coordinates of
a cell are less than the minimum or greater than the maximum specified
coordinate value, then the cell is not marked for refinement.
Similarly, if the cell is located less than the minimum or greater
than the maximum specified distance to a viscous wall, then the cell
is not marked for refinement.  The optional arguments are specified
using the following keywords and corresponding numerical values:
\begin{verbatim}
-xmin <min_x_val>
-xmax <max_x_val>
-ymin <min_y_val>
-ymax <max_y_val>
-zmin <min_z_val>
-zmax <max_z_val>
-dmin <min_d_val>
-dmax <max_d_val>
\end{verbatim}
The coordinate and distance to viscous wall values may be used in any
combination.

\subsection{{\tt marker}: Generating the Refinement Plan}
The tool \verb!marker! is used to produce a refinement plan from the
original mesh and the tag file.  The refinement plan indicates how
cells should be refined (which direction, how many times, {\it{etc.}})
and also keeps track of the original nodes and any newly created
hanging nodes.  The format of the \verb!marker! command is:

\begin{verbatim}
marker -tol <minspace> -fold <fold> -g <orig_vog> <restart_options>
	-o <refplan> -levels <levels> -tag <tag_file>
\end{verbatim}

\noindent
where

\begin{description}
\item {\tt <minspace>} :  smallest mesh size allowed (in meters)
\item {\tt <fold>} :  folding parameter for twisted faces (keep set to 0.25)
\item {\tt <orig\_vog>} :  original \verb!.vog! file name
\item {\tt <restart\_options>} :  none if this run is not a continuation from a previous level; or
\subitem {\tt -restart} :  flag to indicate a continuation from a previous refinement plan (level)
\subitem {\tt -r <rest\_ref\_plan>} :  previous refinement plan file name to use for continuation
\item {\tt <refplan>} :  new refinement plan file name
\item {\tt <levels>} :  number of refinement levels within a refinement step (generally 1, but can be $>$ 1 for 3D isotropic refinement only)
\item {\tt <tag\_file>} :  tag file name
\end{description}

The \verb!-restart! option is important on all mesh refinement levels
except the first ({\it{i.e.}}, the refinement step from the original
mesh to the first refinement level).  The refinement plan keeps track
of the refinement history of each cell and especially how each cell
was formed.  The plan also keeps track of hanging nodes, and all of
this information is needed to maintain mesh quality through subsequent
refinement levels.


\subsection{{\tt refmesh}: Generating the Refined Mesh}

The refinement plan provides the instructions on how to refine each
cell during the current step while \verb!refmesh! carries out the
refinement instructions to generate the new mesh.  The options for
\verb!refmesh! are:

\begin{verbatim}
refmesh -g <orig_vog> -r <refplan> -o <new_vog>
\end{verbatim}

\noindent
where

\begin{description}
\item {\tt <orig\_vog>} :  original \verb!.vog! file name
\item {\tt -r <ref\_plan>} :  refinement plan file name to use
\item {\tt <new\_vog>} :  new (refined) \verb!.vog! file name
\end{description}

\subsection{Curved Boundary and Viscous Mesh Adaptation Tools}

When a highly curved surface mesh is refined, a representation of the
surface geometry is needed in order to allow the newly created mesh
points to remain on the true surface.  The surface geometry is
represented by cubic Bezier patches which can either be created from
the original \verb!.vog! file using \verb!computegeom! or from
SolidMesh using \verb!getgeom!.  Once a mesh is refined, the surface
mesh points are popped back to the original geometry, and the interior
mesh is adjusted to smoothly conform to the geometry using
\verb!adjustpos!.  If a viscous calculation is being refined, hanging
nodes in the highly stretched boundary layer mesh are avoided using
\verb!viscousFix!.  This tool tags points near the viscous boundary
consistently so that if any points within a given delta are tagged for
refinement, then all points within that delta are tagged.


\subsubsection{{\tt computegeom} and {\tt getgeom}: Defining the Boundary Geometry for Surface Mesh Refinement}

This tool is used to define the geometry for surface mesh refinement
of highly curved boundary surfaces.  There are two ways to create the
geometry file (\verb!*.geom!).  The first method is approximate and
uses an existing \verb!.vog! file:
\begin{verbatim}
     computegeom -o sphere.geom sphere
\end{verbatim}
where \verb!sphere.vog! is the original grid file.  The second method is more accurate and uses exact surface normal and geometry information from SolidMesh:
\begin{verbatim}
     getgeom sphere
\end{verbatim}
where \verb!sphere.surf! and \verb!sphere.surfn! are supplied by the user.  The above command then produces \verb!sphere.geom!.


\subsubsection{{\tt viscousFix}: Viscous Mesh Adaptation}

The tool \verb!viscousFix! is used to tag additional points that are
within a specified delta of a viscous boundary.  The code reads the
tag file produced by \verb!adpt! and the topology file produced by
Chem from the \verb!output/! directory.  The user identifies which
boundaries are viscous walls, specifies a delta to use, and inputs the
tag file name and the iteration number.  The format of the command is:
\begin{verbatim}
     viscousFix -bc <viscous_bc_name> -delta <delta_from_wall>
         -tag <tag_file> -o <output_tag_file> <case_name> <iteration>
\end{verbatim}
An example of the command is for 3 viscous walls with a delta of 0.01 at iteration 0 is:
\begin{verbatim}
     viscousFix -bc wall1 -bc wall2 -bc wall3 -delta 0.01
         -tag refine.dat -o vrefine.dat sphere 0
\end{verbatim}
The output tag file \verb!vrefine.dat! should then be used by \verb!marker! to produce the refinement plan.


\subsubsection{{\tt adjustpos}: Preserving the Boundary Geometry}

For the refinement of highly curved boundary surfaces, the linear
refinement pass must be adjusted to allow surface mesh points to pop
back to the boundary geometry and to allow the near surface mesh to
deform smoothly into the domain.  The tool \verb!adjustpos! is used
after a refinement cycle to accomplish this adjustment.  For example,
if we have a linear refined mesh given by \verb!sphere2l.vog!, we can
generate the curvature adjusted mesh \verb!sphere2.vog! using the
command:
\begin{verbatim}
     adjustpos -i sphere2l.vog -o sphere.vog -g sphere.geom
\end{verbatim}
where \verb!sphere.geom! is the geometry file created either by \verb!computegeom! or \verb!getgeom!.

\subsection{Online AMR}

Online adaptation can also be used where the adaption process is
integrated into the CHEM solver.  This can be used for both steady
state and unsteady problems and unlike the command line version,
allows for mesh coarsening as well as refinement.  To enable online
AMR the {\tt adaptFrequency} variable needs to be set in the vars file
which sets the number of timesteps between adaption phases.  A summary
of the online adaption options is as given:

\begin{enumerate}
\item {\tt adaptFrequency}

  How frequently adaption cycles will occur.  If set to a value of
  100, then an adaption cycle will be performed every 100 timesteps.

\item {\tt adaptStart}

  No adaptation will occur before the timestep set by {\tt adaptStart}

\item {\tt adaptSensitivity}

  This sets how mesh elements are selected for refinemnt.  Any cell
  with an error that exceeds the mean error value by this factor times
  the standard deviation of the error will be marked for refinement.
  A value closer to zero will result in more refinement.

\item {\tt coarsenSensitivity}

  This sets how mesh elements will be selected for a coarsening
  operation.  Note that only cells that have been refined can be
  coarsened.  Similar to {\tt adaptSensitivity} this will mark cells
  that are this factor of a standard deviation below the mean error
  for coarsening.

\item {\tt adaptSensor} ({\it list})

  This sets the error sensors that will be utilized for mesh
  refinement.  Currently the sensors are as follows:

  \begin{table}[htbp]

  \begin{centering}
    \leavevmode
  \begin{tabular}{|l|l|}
    \hline
    Variable & description \\
    \hline\hline
    Terror   & Error in the temperature field\\
    Rerror   & Error in the density field \\
    Perror   & Error in the pressure field \\
    OHerror  & Error in OH mass fractions\\
    NOerror  & Error in NO mass fractions\\
    Uerror   & Error in the velocity field\\
    \hline
  \end{tabular}
\caption{Error Estimators for mesh adaption}
  \label{tab:AMRerror}
  \end{centering}
\end{table}

\item {\tt maxAdaptSteps}  ({\it integer})

  The maximum number of adapt cycles that will be performed.

\item {\tt adaptMode} ({\it string})

  The type of adaptation that will be utilized. This can be set to
  'isotropic', 'anisotropic' or '2D'.  Use '2D' for two dimesional
  simulations where the grid is extruded 1 cell in the z-axis.  The
  'isotropic' refinement always splits the cell equally, while the
  'anisotropic' refinement will choose to refine prisms and hexes such
  that the resulting cells tend to have equal mesh spacing in all
  directions.  

\item {\tt adaptMinEdgeLength} ({\it float, units meters})

  Stop refining cells which have a minimum edge length of this size.

\item {\tt adaptGeometryEdgeLength} ({\it float, unit meters})

  Provides a target edge length for refining background grid in
  overset problems.

\item {\tt adaptGeometryLayers} ({\it float, nondimensional})

  Sets thickness of derefinement levels for overset goemetry
  refinement.  This is specified with respect to the {\tt
    adaptGeometryEdgeLength}. 

\item {\tt adaptFaceFold} ({\it float, nondimensional})

  Identifies folded faces to disable refinement of highly twisted
  faces.

\item {\tt adaptBalanceType} ({\it integer})

  Sets method used to balance refinement level between neighboring
  cells.

\item {\tt maxAdaptCellCount} ({\it double})

  Sets the maximum number of cells after refinement specified in
  millions of cells.  Fractional numbers can be used, for example
  setting to $0.1$ would result in a limit of $100,000$ cells.
  
\item {\tt adaptMaskDistance} ({\it double, units meters})

  Sets how close to a solid wall adaptation will be allowed.

\end{enumerate}
